\section{One owner, one table, five states}
\label{sec:one-owner}

The two-$H$ hardness theorem has an exact lower boundary. One $H_a$
already expresses reset, and it retains a genuine owner test. Nevertheless,
its equations have a polynomial decision procedure, uniform in the number
of roles. More generally, the complexity of an individual player's
observation depends on how many \emph{other} owners occur.

Throughout this section the native syntax is
\[
 t::=x_j\mid t\del t\mid H_r(t,t,t),\qquad
 H_r(X,Y,Z)=(Y\star_r Z)\del X.
\]
There are no constants, unions, or internal role switches. Restricting the
$H$ indices does not restrict every node of an owned-choice expansion:
reset itself is a derived operation. Complexity is measured in this native
tree syntax.

\subsection{A root whose status never changes}

Define the root variable by
\[
 \rho(x_j)=j,\qquad \rho(s\del t)=\rho(t),\qquad
 \rho(H_r(x,y,z))=\rho(x).
\]
Write $M_j$ for the assertion that the value of $x_j$ is ambiguous.

\begin{lemma}[Root status]\label{owner:lem:root}
A term is ambiguous exactly when its root variable is ambiguous. For
every role $r$, its value is $\{r\}$ exactly when its root variable is
$\{r\}$.
\end{lemma}
\begin{proof}
Reset returns its right operand when that operand is resolved, and
otherwise returns an ambiguous union containing it. Both status claims
follow by induction, since the right operand in the definition of
$H_r(X,Y,Z)$ is $X$.
\end{proof}

Fix an observer $k$ and one different role $a$. Allow $H_k$ and $H_a$,
in addition to reset. Put
\[
 N_j\ \Longleftrightarrow\ x_j\ne\{a\},\qquad
 V_j\ \Longleftrightarrow\ k\in x_j.
\]
The implications $M_j\Rightarrow N_j$ and $V_j\Rightarrow N_j$ are
the only restrictions needed below. If $t_k$ denotes possible membership
of $k$ in the value of $t$, the exact recursive laws are
\begin{align}
 (s\del t)_k
   &\longleftrightarrow t_k\lor(M_{\rho(t)}\land s_k),\label{owner:eq:reset}\\
 H_k(x,y,z)_k
   &\longleftrightarrow x_k\lor
        (M_{\rho(x)}\land(y_k\lor z_k)),\label{owner:eq:self}\\
 H_a(x,y,z)_k
   &\longleftrightarrow x_k\lor
        (M_{\rho(x)}\land N_{\rho(y)}\land N_{\rho(z)}
                     \land(y_k\lor z_k)).\label{owner:eq:other}
\end{align}
For \eqref{owner:eq:self}, membership of the owner in an owned choice
is disjunction. For \eqref{owner:eq:other}, an outsider survives that
choice precisely when neither child is the absorbing singleton $\{a\}$.
\Cref{owner:lem:root} turns these child tests into input tests.

\subsection{Exact guard comparison}

A guard is a triple $g=(j,A,B)$, with source variable $j$ and finite sets
of input indices $A,B$. Its meaning is
\[
 [g]=V_j\land\bigwedge_{i\in A}M_i\land\bigwedge_{i\in B}N_i.
\]
The list $\Gamma_k(t)$ is built from \eqref{owner:eq:reset}--\eqref{owner:eq:other}:
an input contributes $(j,\varnothing,\varnothing)$; reset retains the
right guards and adds $\rho(t)$ to the $M$ tests of the left guards;
$H$ retains the root guards and adds the displayed tests to each guard
of its other two arguments. Consequently
\begin{equation}\label{owner:eq:guards}
 t_k\longleftrightarrow\bigvee_{g\in\Gamma_k(t)}[g].
\end{equation}
Each variable occurrence contributes exactly one guard. No distributive
expansion duplicates a guard.
For example, the outsider guards of $H_a(x,y,z)$ are
\[
 (x,\varnothing,\varnothing),\qquad
 (y,\{x\},\{y,z\}),\qquad (z,\{x\},\{y,z\}).
\]
In the second guard the $N_y$ test is automatic from its source $V_y$;
similarly $N_z$ is automatic in the third.

Say $h=(i,C,D)$ \emph{covers} $g=(j,A,B)$ when
\begin{equation}\label{owner:eq:cover}
 i=j,\qquad C\subseteq A,\qquad D\subseteq A\cup B\cup\{j\}.
\end{equation}
These conditions imply $[g]\Rightarrow[h]$; the last extra index $j$
uses $V_j\Rightarrow N_j$.

\begin{theorem}[Guard completeness]\label{owner:thm:guards}
Suppose $|I|\geq3$, $k\ne a$, and all $H$ indices in $s,t$ belong to
$\{k,a\}$. Then $s_k\Rightarrow t_k$ on every nonempty valuation if
and only if every guard of $\Gamma_k(s)$ is covered by a guard of
$\Gamma_k(t)$. A failed comparison has an explicitly constructible
singleton-or-pair countervaluation using just three roles.
\end{theorem}
\begin{proof}
Sufficiency follows from \eqref{owner:eq:guards} and
\eqref{owner:eq:cover}. For necessity, choose $b\notin\{k,a\}$ and
an uncovered guard $g=(j,A,B)$. Assign
\[
 x_i=\begin{cases}
 \{k,a\},&i=j\text{ and }i\in A,\\
 \{k\},&i=j\text{ and }i\notin A,\\
 \{a,b\},&i\ne j\text{ and }i\in A,\\
 \{b\},&i\ne j,\ i\notin A,\ i\in B,\\
 \{a\},&\text{otherwise}.
 \end{cases}
\]
Exactly the indices in $A$ are ambiguous; exactly those in
$A\cup B\cup\{j\}$ satisfy $N$; only input $j$ contains $k$.
Thus $g$ holds, and a right guard holds exactly when it covers $g$.
No right guard does. This is the required countervaluation.
\end{proof}

This also gives canonical tables: for each source, remove its index and
the $M$ indices from its $N$ tests, then retain only the minimal guards
under \eqref{owner:eq:cover}. Minimal antichains are unique because
\cref{owner:thm:guards} characterizes implication exactly.
For trees of total size $n$, there are at most $n$ guards, each with
$O(n)$ tests. Direct list-based compilation and all-pairs coverage give
an $O(n^4)$ equality bound when input names have constant-time equality;
bit-set implementations improve the elementary subset checks. The
algorithm is uniform in $K$ and never enumerates winner sets.
An affirmative comparison has a short certificate: provide, for each
left guard, the index of a covering right guard. Guard compilation and
the three tests in \eqref{owner:eq:cover} check the certificate directly.
This is a complete comparison calculus uniform in $K$; it does not
assert a finite equational basis for the native $H$ signature.

\subsection{A single outsider sees every equation}

Now allow just one $H$ index $a$. All outsider tables have the same
syntax. The owner's table is simpler still:
\begin{equation}\label{owner:eq:forget}
 \Gamma_a(t)=\{(j,A,\varnothing):(j,A,B)\in\Gamma_k(t)\},
 \qquad k\ne a.
\end{equation}
This identity is for the unminimized guard lists. It follows directly
from the recursion: an owner's own $H$ adds no $N$ tests. Forgetting
the $N$ tests preserves coverage.

\begin{theorem}[One-table completeness]\label{owner:thm:one-table}
For $K\geq3$, terms in $\{\del,H_a\}$ are equal on every nonempty
valuation if and only if their guard lists for any one outsider cover
each other. Equation validity is decidable in polynomial time.
\end{theorem}
\begin{proof}
Equality implies observational equality and hence mutual coverage by
\cref{owner:thm:guards}. Conversely, all outsiders have the same lists,
so mutual coverage gives equality of every outsider's membership.
By \eqref{owner:eq:forget}, it also gives mutual coverage for the owner,
hence equality of that membership. These memberships determine the set.
\end{proof}

There is a small algebra behind this result. For a fixed outsider $k$,
retain the three flags
\[
 q(S)=(m,e,p)=
   (|S|>1,\ S=\{a\},\ k\in S).
\]
There are exactly five possible states:
\begin{center}
\begin{tabular}{lccc l}
\toprule
State & $m$ & $e$ & $p$ & Three-role representative\\
\midrule
Owner singleton & 0&1&0&$\{a\}$\\
Observed singleton & 0&0&1&$\{k\}$\\
Other singleton & 0&0&0&$\{b\}$\\
Observed ambiguity & 1&0&1&$\{k,a\}$\\
Other ambiguity & 1&0&0&$\{a,b\}$\\
\bottomrule
\end{tabular}
\end{center}
Write $d(p)=(1,0,p)$. On these five states define
\begin{align*}
 u\del v&=\begin{cases}v,&\neg m_v,\\
                         d(p_u\lor p_v),&m_v,\end{cases}\\
 h(u,v,w)&=\begin{cases}u,&\neg m_u,\\
 d(p_u\lor(\neg e_v\land\neg e_w\land(p_v\lor p_w))),&m_u.
 \end{cases}
\end{align*}
Call this algebra $\mathbf F$.

\begin{theorem}[A common five-element generator]\label{owner:thm:five}
For every $K\geq3$ and role $a$, the algebra
$(\Pow(I);\del,H_a)$ has exactly the identities of $\mathbf F$.
In particular its equational theory is independent of $K$.
\end{theorem}
\begin{proof}
The displayed recursion makes $q$ a surjective homomorphism; the five
representatives prove surjectivity. Thus every identity of the winner-set
algebra holds in $\mathbf F$. Conversely, an identity holding in
$\mathbf F$ gives equality of the outsider flag $p$ on every valuation.
By \cref{owner:thm:guards}, this implies mutual guard coverage, and
\cref{owner:thm:one-table} supplies equality of the complete winner sets.
\end{proof}

The outsider suffices here at the level of identities: two terms with
the same outsider observation \emph{on all valuations} are equal.
This does not say that the five-state map is injective on winner sets,
or that one observer determines the value of an individual execution.

\subsection{The exact owner boundary}

\begin{corollary}[Owner-count dichotomy]\label{owner:cor:complexity}
Fix $K\geq3$ and $R\subseteq I$. For the native signature consisting of
reset and $H_r$ for $r\in R$, full equation validity is polynomial-time
decidable when $|R|\leq1$, and coNP-complete when $|R|\geq2$.
For a fixed observer $k$, observational implication and equivalence are
polynomial-time decidable when $|R\setminus\{k\}|\leq1$, and
coNP-complete when $|R\setminus\{k\}|\geq2$.
\end{corollary}
\begin{proof}
\Cref{owner:thm:one-table,owner:thm:guards} give the algorithms; unused
owner indices may be added to apply them. Two distinct $H$ operations
give the full-equation lower bound by \cref{thm:H-hardness}.
For the observational lower bound choose two owners $a,b\ne k$ in that
same reduction and use $c=\{k\}$ as its fresh-marker witness. In an
unsatisfiable instance the equation is valid for all roles; in a
satisfiable instance this valuation separates membership of $k$.
The reduction is to the implication $c\del R\subseteq_k R$, where
$\subseteq_k$ means implication of $k$ membership, as well as to its
equivalence. Nonvalidity has a polynomial singleton-or-pair certificate,
so both problems belong to coNP.
\end{proof}

Thus with $H_a,H_b$, both participating owners have polynomial
observational equality, while a third player's observation is already
coNP-complete. Even agreement of both owners is insufficient. Put
\[
 W=y_1\del(y_2\del z),\quad
 T=H_a(z,y_1,c)\del(H_b(z,y_2,c)\del W).
\]
The terms $c\del T$ and $T$ have identical possible-$a$ and possible-$b$
observations on every valuation. Indeed, a potential new $a$ marker
already in $W$ needs no argument; one outside $W$ has $y_1\ne\{a\}$,
and survives through $H_a(z,y_1,c)$. The same argument uses $y_2$ for
$b$. A resolved $z$ makes both terms equal to $z$.
But with $y_1=\{a\}$, $y_2=\{b\}$, $z=\{a,b\}$, $c=\{k\}$,
the values are $\{a,b,k\}$ and $\{a,b\}$ respectively.

The source of tractability is the preserved root status and the single
chain of outsider tests, not atom preservation by itself. Union loses
that root law and already gives owner-free hardness. A second other
owner supplies incompatible singleton tests and recovers Boolean
constraints. The guard criterion, executable comparison, one-table
completeness, five-element theorem, and compilation to the existing
owned-choice semantics are kernel-checked in
\path{lean/OwnerFragments.lean}. The polynomial size bounds and the
complexity deductions are written arguments.
